\begin{resumen}%
En esta tesis estudiamos el movimiento de la tortuga terrestre argentina, \textit{Chelonoidis chilensis}, una especie
en peligro y poco estudiada, a partir de un año de registros de GPS de seis individuos de una población de Río Negro.
Entender cómo se mueve y cómo usa su hábitat aporta bases cuantitativas para su conservación.

Primero caracterizamos el error del GPS: la diferencia entre posiciones de una tortuga quieta sigue una distribución
t de Student de colas pesadas, y el error tiene memoria en dos escalas, de unos 30 minutos y de algunas horas. Para un
animal tan lento, este error es comparable al movimiento y debe incluirse en el modelo.

Luego construimos un modelo estocástico de \textit{Stop and Go}, con duraciones de caminata y reposo en ley de potencias
tomadas de mediciones previas con acelerómetros, y le agregamos por capas el error del GPS, una hora de salida del
refugio sigmoide y una velocidad variable entre días. Este modelo reproduce la distribución de desplazamientos a 15
minutos y el desplazamiento cuadrático medio diario, pero no los tipos de día observados. Clasificamos los días en
reposo, exploración y retorno, y cambio de refugio, y modelamos cada tipo por separado: combinados, los tres submodelos
reproducen todas las métricas analizadas.

Finalmente, describimos cómo cambia el comportamiento a lo largo del año: en la reproducción las tortugas usan más
espacio, cambian más de refugio y salen más temprano que en la brumación. La selección de refugios sigue una dinámica
de exploración y retorno preferencial similar a la observada en seres humanos.

\end{resumen}

\begin{abstract}%
\end{abstract}




%%% Local Variables: 
%%% mode: latex
%%% TeX-master: "template"
%%% End: 
